
\chapter*{Présentation de l’IP2I}
\justifying
\sloppy

L’\textbf{Institut de Physique des 2 Infinis de Lyon (IP2I)} est une unité mixte de recherche placée sous la tutelle du \textbf{CNRS} et de l’\textbf{Université Claude Bernard Lyon 1}. L’institut développe des activités de recherche couvrant notamment la physique des particules, la physique nucléaire, l’astrophysique, la cosmologie, ainsi que les applications des rayonnements ionisants. Cette diversité scientifique s’accompagne d’un fort volet instrumental, au sein duquel la conception de détecteurs, de systèmes électroniques et de chaînes de lecture occupe une place structurante.

L’IP2I se distingue par un fonctionnement fortement interdisciplinaire. Chercheurs, enseignants-chercheurs, ingénieurs, techniciens et doctorants y collaborent autour de projets dans lesquels la qualité des dispositifs expérimentaux conditionne directement la pertinence des résultats scientifiques. Dans ce cadre, l’ingénierie n’intervient pas comme un simple support, mais comme une composante essentielle de la démarche de recherche, depuis l’expression du besoin jusqu’à la validation des solutions mises en œuvre.

\subsection{\textbf{Présentation de l’équipe de microélectronique}}

Au sein de cet environnement, l’équipe de \textbf{microélectronique} contribue au développement de circuits intégrés dédiés à l’instrumentation scientifique, en particulier pour les détecteurs nécessitant des performances élevées en précision temporelle, en compacité et en robustesse. Son activité couvre l’analyse du besoin, la définition d’architectures, la conception de blocs analogiques, mixtes et numériques, la simulation, la validation et la préparation à l’intégration dans des systèmes plus complets.

L’organisation de l’équipe repose sur deux pôles complémentaires :
\begin{itemize}
    \item \textbf{le pôle conception analogique et mixte}, chargé de l’étude et de la conception des blocs de front-end, des circuits de polarisation, des fonctions de synchronisation et, plus largement, des éléments analogiques ou mixtes nécessaires au traitement fin du signal ;
    \item \textbf{le pôle conception et vérification numérique}, dédié à la définition des architectures numériques, au développement RTL, à la vérification fonctionnelle et à l’intégration des interfaces de lecture, en s’appuyant notamment sur des méthodologies telles que le \textit{Digital on Top} (DoT) et l’\textit{Universal Verification Methodology} (UVM).
\end{itemize}

Cette structuration favorise une approche transversale des projets. Les interactions entre analogique, mixte et numérique sont en effet constantes dès lors qu’il s’agit de concevoir des \textit{ASIC} de détection à haute performance. L’équipe de microélectronique collabore ainsi étroitement avec les physiciens expérimentaux, qui définissent les besoins scientifiques, mais aussi avec les autres groupes techniques impliqués dans la lecture, l’acquisition et l’exploitation des données.

\subsection{\textbf{Inscription du travail réalisé dans l’équipe}}

Le travail présenté dans ce rapport s’inscrit pleinement dans cette dynamique collective. Il porte sur l’étude et l’implémentation d’un bloc de \textbf{conversion temps-numérique} destiné à être intégré dans le chip \textbf{SPADMIC}. Ce sujet se situe à l’interface entre plusieurs problématiques : compréhension théorique de la mesure temporelle fine, choix d’architecture, organisation de la lecture des données et prise en compte des contraintes d’intégration dans un système complet.

Dans ce contexte, l’environnement offert par l’IP2I et par son équipe de microélectronique constitue un cadre particulièrement pertinent. Il permet d’aborder le sujet avec une vision à la fois scientifique, technologique et méthodologique, en reliant les besoins de la détection aux choix de conception retenus pour l’architecture étudiée dans la suite de ce rapport.

\clearpage
